\section{Problem setting}
\label{sec:problem}

We consider a turbulent flow in an urban geometry that is forced
through a small set of inflow parameters
$\param(t) \in \R^{n_p}$ assumed to vary in time.  In our experiments
$\param$ contains the inflow angle and the inflow speed, optionally
augmented by a bulk pressure-gradient magnitude, so $n_p \in \{2,
3\}$.  The flow state $\state(t)$ -- velocity components together
with any auxiliary fields produced by the solver -- evolves under a
forward model
\begin{equation}
  \state(t) = \fwd\bigl(\param(\cdot)\big|_{[0,t]};\, \state(0)\bigr),
  \label{eq:forward-abstract}
\end{equation}
which is in general non-linear, expensive, and only available as a
non-differentiable black box.  In this work $\fwd$ is one of three
urban-flow solvers (\pylbm{}, \pyudales{}, \pypalm{}) wrapped behind
a common Python interface (Section~\ref{sec:forward}).

\paragraph{Observations.}
A linear observation operator $\obsop$ extracts pointwise velocity
samples at a small set of probe locations and aggregates them in
time over fixed-length intervals.  Observations are assumed to be
contaminated by additive zero-mean Gaussian noise,
\begin{equation}
  \obs = \obsop(\state) + \noise, \qquad \noise \sim \mathcal{N}(\zerovec, \CD),
  \label{eq:obs-model}
\end{equation}
with diagonal covariance $\CD = \sigma_d^2 \identity$.  Throughout
this report we work in an identical-twin
setting: the truth is generated by the forward model with a known
parameter trajectory $\paramTraj^\star = \{\param^\star(t)\}_t$, and
synthetic observations are produced by sampling~\eqref{eq:obs-model}
along the truth.  To avoid the inverse crime of using identical
generative processes for truth and prior, the truth trajectory is
drawn from a different parameter time-series generator (typically
with a different correlation length) than the assimilation prior.

\paragraph{Time discretization.}
The full simulation horizon $[0, \Nw\, \windowLen]$ is partitioned
into $\Nw$ contiguous windows of length $\windowLen$.  Within each
window the parameter trajectory is represented on $\Nt$ discrete
time points, and adjacent windows share their boundary point so that
the underlying truth and any extrapolation between windows are
continuous in time.  Each window's simulation is run with its own
\emph{local} clock $[0, \windowLen]$ and warm-starts from the
previous window's terminal state; absolute times are reconstructed
only at save and plot time from the window index.

\paragraph{Inverse problem.}
Given noisy observations $\obs\win{w}$ collected during window $w$
and a prior over the parameter trajectory of that window, the goal
is to characterize the posterior
\begin{equation}
  p\bigl(\paramTraj\win{w} \mid \obs^{(0)}, \ldots, \obs\win{w}\bigr)
  \label{eq:posterior}
\end{equation}
through an ensemble of $\Ne$ approximate samples.  Because direct
sampling from~\eqref{eq:posterior} is infeasible -- both because of
the cost of $\fwd$ and because of the non-Gaussianity introduced by
the observation operator and the time-series prior -- we approximate
each conditional update by a tempered ensemble Kalman analysis
(Section~\ref{subsec:da-single-window}) and propagate the dependence
on past observations forward in time through the parameter
time-series model (Section~\ref{subsec:da-prior}).
